\section{Agent Architecture \& Directory Structure}

\subsection{Repository Organization}
The repository is separated into requirement sources, reusable project knowledge, agent memory, generated deliverables, execution logs, and OpenCode configuration. Each directory has a specific role, which prevents the agent from mixing source material with final outputs or temporary evaluation artifacts. Table~\ref{tab:repo-structure} maps the major paths to their responsibilities.

\begin{table}[htbp]
\centering
\footnotesize
\renewcommand{\arraystretch}{1.18}
\begin{tabularx}{\textwidth}{>{\raggedright\arraybackslash}p{4.0cm} X}
\toprule
\textbf{Path} & \textbf{Responsibility} \\
\midrule
\texttt{README.md} & Provides setup instructions, run guidance, and the list of required outputs. \\
\texttt{AGENTS.md} & Defines agent rules and project-specific operating instructions. \\
\texttt{CS486\_Project.pdf} & Stores the original assignment brief used as the external project reference. \\
\texttt{req/} & Contains the source business requirements, especially \texttt{business-requirement.md}. \\
\texttt{docs/} & Acts as the project knowledge base and workflow constraints, including the overview, tech stack, entity registry, schema registry, and design decisions. \\
\texttt{docs/templates/} & Provides templates and reading or writing guidance for registry-style documentation. \\
\texttt{memory/} & Stores persistent state such as \texttt{MEMORY.md}, \texttt{ActiveContext.md}, and \texttt{Progress.md}. \\
\texttt{outputs/} & Contains final deliverables for Tasks 01--07, from requirement analysis to business query design. \\
\texttt{logs/trajectory/} & Stores task execution traces, grouped by task number. \\
\texttt{logs/eval/} & Stores output evaluation records and future pipeline-level evaluation logs. \\
\texttt{.opencode/commands/} & Defines slash commands for running the full database design pipeline, generating individual tasks, and evaluating results. \\
\texttt{.opencode/skills/} & Stores the main database design pipeline skill, task-specific skills, and the evaluation framework. \\
\bottomrule
\end{tabularx}
\caption{Major repository paths and their responsibilities}
\label{tab:repo-structure}
\end{table}



\subsection{Layered Documentation Strategy}
The project adopts an AI-first documentation strategy based on the principle of progressive disclosure~\cite{viblo-ai-agent-convention}. Rather than loading the entire repository into the model context at the beginning of each session, information is organized into multiple layers that are accessed only when relevant to the current task. This design reduces context overload, improves retrieval efficiency, and ensures that stable project knowledge is separated from task-specific procedures and execution history. Table~\ref{tab:doc-layers} summarizes the four documentation layers used by the agent.



\begin{table}[htbp]
\centering
\footnotesize
\renewcommand{\arraystretch}{1.18}
\begin{tabularx}{\textwidth}{>{\raggedright\arraybackslash}p{0.8cm}
                            >{\raggedright\arraybackslash}p{4.7cm}
                            >{\raggedright\arraybackslash}p{3.7cm}
                            >{\raggedright\arraybackslash}X}
\toprule
\textbf{Layer} & \textbf{Primary files} & \textbf{When the agent reads them} & \textbf{Purpose} \\
\midrule
1 & \texttt{AGENTS.md} & At the start of the session & Defines global rules, editing constraints, and workflow expectations. \\
2 & \texttt{docs/README.md}, \texttt{docs/*.md}, \texttt{docs/templates/} & On demand, following the documented reading order & Stores domain knowledge, naming conventions, design decisions, entity definitions, schema definitions, and reusable templates. \\
3 & \texttt{.opencode/commands/}, \texttt{.opencode/skills/} & When a command or task workflow is invoked & Supplies command entry points, pipeline procedures, task rules, and evaluation guidance. \\
4 & \texttt{memory/*.md}  & During and after task execution & Preserves active context, progress, corrections. \\
\bottomrule
\end{tabularx}
\caption{Progressive documentation layers used by the agent workflow}
\label{tab:doc-layers}
\end{table}



\subsection{Agent Workflow}

The agent architecture is implemented as a controlled database-design pipeline rather than a single free-form prompting process. When a user invokes a task, the request is first routed to an OpenCode command, which acts as the workflow entry point. The command identifies the requested stage of the database design lifecycle and loads the corresponding task-specific skill. Examples include business requirement analysis, conceptual database design, logical schema design, database validation, SQL implementation, sample data generation, query design, and evaluation.

Unlike a traditional prompt-based workflow, the system does not rely solely on conversation history to preserve context. Instead, important design decisions are externalized into persistent knowledge artifacts. During execution, the selected skill consults the project knowledge base, reusable templates, and previously generated registries before producing an output. As a result, information discovered in earlier tasks can be reused consistently throughout the pipeline without requiring the model to rediscover the same knowledge multiple times.

The workflow follows a clear separation of concerns. OpenCode commands define task entry points, skills define generation procedures and evaluation rules, the \texttt{docs/} directory serves as the authoritative knowledge base, the \texttt{memory/} directory maintains persistent project state, the \texttt{outputs/} directory stores final deliverables, and the \texttt{logs/} directory records execution traces and evaluation evidence. This modular organization improves reproducibility, simplifies auditing, and reduces unintended coupling between different stages of the workflow.

The architecture also incorporates a feedback mechanism. After each task is completed, trajectory logs and evaluation reports are generated and reviewed. Important corrections, unresolved issues, and validated design decisions can then be propagated back into memory and project documentation. Consequently, future tasks benefit from accumulated project knowledge while remaining independent of the original conversation context. Figure~\ref{fig:agent-workflow} illustrates the overall workflow.


\begin{figure}[h]
\centering
\resizebox{\textwidth}{!}{%
\begin{tikzpicture}[
    node distance=1.05cm and 1.15cm,
    box/.style={rectangle, rounded corners, draw, minimum width=2.75cm, minimum height=0.85cm, align=center, font=\small},
    smallbox/.style={rectangle, rounded corners, draw, minimum width=2.35cm, minimum height=0.75cm, align=center, font=\footnotesize},
    widebox/.style={rectangle, rounded corners, draw, minimum width=3.25cm, minimum height=0.9cm, align=center, font=\small},
    arrow/.style={-{Stealth}, thick},
    feedback/.style={-{Stealth}, thick, dashed},
]

% Main workflow row
\node[box] (user) {User Request};
\node[box, right=of user] (onboard) {Session\\Onboarding};
\node[box, right=of onboard] (select) {Identify\\Active Task};
\node[widebox, right=of select] (load) {Load Command,\\Skill, Template};
\node[widebox, right=of load] (generate) {Generate or\\Revise Artifact};

% Lower row
\node[box, below=of select] (pipeline) {Selected Stage\\Task 01--07};
\node[box, below=of load] (validate) {Validate Against\\Rules \& Registries};
\node[box, below=of generate] (logs) {Write Evidence\\Logs};
\node[box, right=of logs] (review) {User Review\\Approval Gate};

% Context inputs
\node[smallbox, above=of onboard] (agents) {Operating Rules\\\texttt{AGENTS.md}};
\node[smallbox, above=of select] (memory) {Project State\\\texttt{memory/}};
\node[smallbox, above=of load] (docs) {Knowledge Base\\\texttt{docs/}};
\node[smallbox, above=of generate] (req) {Requirements\\\texttt{req/}};

% Outputs / persistent records
\node[smallbox, below=of validate] (outputs) {Deliverables\\\texttt{outputs/}};
\node[smallbox, below=of logs] (trajectory) {Trajectory / Eval\\\texttt{logs/}};
\node[smallbox, below=of review] (progress) {Approved Updates\\\texttt{Progress.md}\\\texttt{ActiveContext.md}};

% Main arrows
\draw[arrow] (user) -- (onboard);
\draw[arrow] (onboard) -- (select);
\draw[arrow] (select) -- (load);
\draw[arrow] (load) -- (generate);
\draw[arrow] (generate) -- (logs);
\draw[arrow] (logs) -- (review);

% Context arrows
\draw[arrow] (agents) -- (onboard);
\draw[arrow] (memory) -- (select);
\draw[arrow] (docs) -- (load);
\draw[arrow] (req) -- (generate);
\draw[arrow] (select) -- (pipeline);
\draw[arrow] (pipeline) -- (load);

% Validation and storage arrows
\draw[arrow] (generate) -- (validate);
\draw[arrow] (validate) -- (outputs);
\draw[arrow] (validate) -- (logs);
\draw[arrow] (logs) -- (trajectory);

% Approval and feedback
\draw[arrow] (review) -- (progress);
\draw[feedback] (review.north) .. controls +(0,1.2) and +(0,-1.2) .. node[above, font=\footnotesize] {revision request} (load.south);


% Bounding boxes
\begin{pgfonlayer}{background}
    \node[draw=gray, dashed, rounded corners,
          fit=(agents)(memory)(docs)(req),
          label=above:{\small Context Sources}] {};
    \node[draw=gray, dashed, rounded corners,
          fit=(pipeline)(load)(generate)(validate),
          label=below:{\small Task Execution Layer}] {};
    \node[draw=gray, dashed, rounded corners,
          fit=(outputs)(trajectory)(progress),
          label=below:{\small Persistent Results Layer}] {};
\end{pgfonlayer}

\end{tikzpicture}%
}
\caption{Agent Workflow}
\label{fig:agent-workflow}
\end{figure}



\subsection{Knowledge Propagation Across Tasks}
A key design principle of the pipeline is knowledge propagation. Rather than treating each task as an isolated generation step, the system stores important outputs as reusable artifacts that become inputs for subsequent tasks. This approach improves consistency across the database design lifecycle and reduces the risk of contradictory decisions between stages.Knowledge propagation of these artifacts are detailed in Table~\ref{tab:task-artifacts}.

\begin{table}[htbp]
\centering
\scriptsize
\renewcommand{\arraystretch}{1.22}
\setlength{\tabcolsep}{3pt}
\begin{tabularx}{\textwidth}{>{\centering\arraybackslash}p{0.75cm}
                            >{\raggedright\arraybackslash}p{3.25cm}
                            >{\raggedright\arraybackslash}p{3.65cm}
                            >{\raggedright\arraybackslash}X}
\toprule
\textbf{Task} & \textbf{Main input(s)} & \textbf{Main output(s)} & \textbf{Role in the pipeline} \\
\midrule
01 & \texttt{req/business-}\linebreak\texttt{requirement.md} & \texttt{outputs/01-business-}\linebreak\texttt{req-analysis-G05.md}\linebreak\texttt{docs/entity-}\linebreak\texttt{registry.md} & Extracts the business purpose, actors, business rules, candidate entities, attributes, and relationships from the source requirements. \\

02 & \texttt{docs/entity-}\linebreak\texttt{registry.md}\linebreak\texttt{outputs/01-business-}\linebreak\texttt{req-analysis-G05.md} & \texttt{outputs/02-erd-}\linebreak\texttt{design-G05.md}\linebreak\texttt{updated entity registry} & Converts conceptual entities into an ERD, confirms relationships, cardinalities, participation constraints, and junction entities. \\

03 & \texttt{docs/entity-}\linebreak\texttt{registry.md}\linebreak\texttt{outputs/02-erd-}\linebreak\texttt{design-G05.md} & \texttt{outputs/03-logical-}\linebreak\texttt{design-G05.md}\linebreak\texttt{docs/schema-}\linebreak\texttt{registry.md} & Translates the ERD into a normalized relational schema, finalizes names and data types, and populates the schema registry. \\

04 & \texttt{docs/entity-}\linebreak\texttt{registry.md}\linebreak\texttt{docs/schema-}\linebreak\texttt{registry.md}\linebreak\texttt{outputs/03-logical-}\linebreak\texttt{design-G05.md} & \texttt{outputs/04-design-}\linebreak\texttt{validation-G05.md}\linebreak\texttt{schema-freeze decision} & Validates requirement coverage, ERD fidelity, normalization, and schema consistency before freezing the schema. \\

05 & \texttt{docs/schema-}\linebreak\texttt{registry.md}\linebreak\texttt{outputs/04-design-}\linebreak\texttt{validation-G05.md} & \texttt{outputs/05-db-}\linebreak\texttt{definition-G05.sql}\linebreak\texttt{DDL compile logs} & Generates SQL Server DDL from the frozen schema, including tables, constraints, indexes, triggers, and validation-supporting database logic. \\

06 & \texttt{outputs/05-db-}\linebreak\texttt{definition-G05.sql}\linebreak\texttt{docs/schema-}\linebreak\texttt{registry.md}\linebreak\texttt{docs/entity-}\linebreak\texttt{registry.md} & \texttt{outputs/06-sample-}\linebreak\texttt{data-G05.sql}\linebreak\texttt{logs/execution/task06/}\linebreak\texttt{logs/eval/task06/} & Produces realistic, idempotent sample data and expected-error tests that prove key business rules against the approved schema. 
\\
07 & \texttt{outputs/05-db-}%
\linebreak\texttt{definition-G05.sql}%
\linebreak\texttt{outputs/06-sample-}%
\linebreak\texttt{data-G05.sql}%
\linebreak\texttt{outputs/01-business-}%
\linebreak\texttt{req-analysis-G05.md}%
\linebreak\texttt{req/business-}%
\linebreak\texttt{requirement.md} & \texttt{outputs/07-query-}%
\linebreak\texttt{design-G05.sql}%
\linebreak\texttt{logs/eval/task07/}%
 &  To produce personalized, meaningful, parameterized, and validated T-SQL queries based on student ideas that directly resolve core operational and analytic questions from role-based perspectives without loss of prior work.   \\

\bottomrule
\end{tabularx}
\caption{Task-specific artifacts used across the database design pipeline}
\label{tab:task-artifacts}
\end{table}